
\chapter{導入}\label{chap:introduction}

% ============================================================
\section{最適輸送とは何か}\label{sec:what-is-ot}

最適輸送 (Optimal Transport; OT) の考え方は，
身近な例を通じて直感的に理解できる．

\begin{example}{砂山の移動}{sand}
ある場所に砂山がある．
この砂をすべて使って，別の場所に指定された形の砂山を作りたい．
砂を運ぶにはコスト（距離や労力）がかかる．
\textbf{総コストを最小にするように砂を移す方法}を求める問題が，
最適輸送問題の原型である．
\end{example}

この問題を数学的に言い換えると，
「ある確率分布 $\mu$ を別の確率分布 $\nu$ に変換する
最もコストの低い方法を見つけよ」となる．

\medskip

一つの物を最短経路で運ぶのは容易であるが，
\textbf{分布全体を同時に移す}場合は質的に異なる難しさがある．
分布の各部分がどこへ向かうかを同時に決定しなければならず，
全体としての最適性を考える必要があるからである．

この問題の考察は18世紀の Monge~\cite{Monge1781} にまで遡り，
以来 200 年以上にわたって数学の中心的な問題の一つであり続けてきた．
近年では，確率分布間の「距離」としての側面や，
効率的な計算手法の発展を通じて，
機械学習やコンピュータビジョンの分野でも中核的なツールとなっている．

% --- 簡単な図 ---
\begin{figure}[htbp]
\centering
\begin{tikzpicture}[>=Stealth, scale=0.9]
  % 左の分布
  \foreach \x/\h in {0/1.2, 0.6/0.8, 1.2/1.6, 1.8/0.4} {
    \fill[black!25] (\x, 0) rectangle (\x+0.5, \h);
  }
  \node[below] at (1.1, -0.3) {分布 $\mu$};

  % 右の分布
  \foreach \x/\h in {5/0.6, 5.6/1.4, 6.2/1.0, 6.8/1.0} {
    \fill[black!50] (\x, 0) rectangle (\x+0.5, \h);
  }
  \node[below] at (6.1, -0.3) {分布 $\nu$};

  % 矢印
  \draw[->, thick, dashed] (2.6, 0.8) -- (4.7, 0.8)
    node[midway, above] {\small 最小コストで輸送};
\end{tikzpicture}
\caption{最適輸送問題の概要：
分布 $\mu$ の「質量」を分布 $\nu$ へ，総コスト最小で移す．}
\label{fig:ot-overview}
\end{figure}


% ============================================================
\section{歴史的背景}\label{sec:history}

最適輸送の研究は，異なる時代・異なる分野で独立に発展してきた．
以下にその流れを整理する．

\subsection*{Monge の定式化（1781年）}

フランスの数学者 Gaspard Monge は，
土木工事における土砂の運搬問題として最適輸送を定式化した~\cite{Monge1781}．
彼の問題は「各粒子をどこに送るか」を写像 $T$ で記述するもので，
これが後に \textbf{Monge 問題}と呼ばれる．

Monge の定式化は直感的だが，
写像 $T$ の存在が保証されない場合があるなど，
数学的に扱いづらい側面があった．

\subsection*{ロジスティクスと経済学（1930--40年代）}

Tolstoi~\cite{Tolstoi1930}（1930年）や
Hitchcock~\cite{Hitchcock1941}（1941年）は，
貨物輸送の最適化という実務的な動機から同種の問題を研究した．

Kantorovich~\cite{Kantorovich1942}（1942年）は，
Monge の写像 $T$ を\textbf{輸送計画}（カップリング）$\pi$
に緩和するという画期的なアイデアを導入した．
この緩和により問題は線形計画として定式化でき，
双対理論を通じた豊かな数学的構造が明らかになった．
Kantorovich はこの業績により1975年にノーベル経済学賞を受賞している．

\subsection*{線形計画法との接続（1949年--）}

Dantzig~\cite{Dantzig1951} が開発した単体法は，
Kantorovich 緩和を数値的に解く最初の汎用的な手段を与えた．
これにより最適輸送は，理論的な対象であると同時に
\textbf{計算可能な}最適化問題となった．

\subsection*{純粋数学への回帰（1990年代）}

1990年代に入り，
Brenier~\cite{Brenier1991} が凸性や偏微分方程式との
深い関係を明らかにしたことで，
最適輸送は純粋数学において大きな復興を遂げた．
Villani による二つの教科書~\cite{Villani2003, Villani2009}
がこの理論の体系化に大きく貢献している．

\subsection*{計算手法の革新と応用の拡大（2000年代--）}

2000年頃から，コンピュータビジョンの分野で
最適輸送コストが \textbf{Earth Mover's Distance (EMD)}
の名で注目を集めた~\cite{Rubner2000}．
しかし，$n$ 点の離散分布に対する厳密解の計算量は $O(n^3 \log n)$ であり，
大規模データへの適用は困難であった．

この状況を劇的に変えたのが，
Cuturi~\cite{Cuturi2013} による
\textbf{エントロピー正則化}と \textbf{Sinkhorn アルゴリズム}の提案である．
行列のスケーリング操作の反復のみで近似解が得られるため，
GPU 上での並列計算とも相性がよく，
計算量は $O(n^2)$ 程度にまで削減される．
これを契機として，最適輸送は
画像処理，自然言語処理，生成モデル，グラフ学習など
データサイエンスの広範な領域で実用的なツールとなった．

% --- タイムライン ---
\begin{figure}[htbp]
\centering
\begin{tikzpicture}[>=Stealth, scale=0.85]
  % 時間軸
  \draw[->, thick] (-0.5, 0) -- (15.5, 0);
  % 目盛りとラベル
  \foreach \x/\yr/\lab in {
    0/1781/{Monge},
    2.5/1930/{Tolstoi},
    3.8/1942/{Kantorovich},
    5.5/1949/{Dantzig (LP)},
    8/1991/{Brenier},
    10.5/2000/{EMD},
    13/2013/{Sinkhorn}%
  } {
    \fill (\x, 0) circle (2.5pt);
    \node[below, font=\footnotesize] at (\x, -0.15) {\yr};
    \node[above, font=\footnotesize, text width=2cm, align=center]
      at (\x, 0.15) {\lab};
  }
\end{tikzpicture}
\caption{最適輸送研究の主な転換点．}
\label{fig:timeline}
\end{figure}


% ============================================================
\section{なぜ最適輸送が重要か}\label{sec:why-ot}

最適輸送が多くの分野で注目される理由は，
それが確率分布間の「距離」として
他の指標にはない優れた性質を持つことにある．

\paragraph{(1) 台の幾何を反映する．}
KL ダイバージェンスなどの $f$-ダイバージェンスは
密度関数の比のみに基づくため，
データが「空間のどこにあるか」という情報を直接には使わない．
一方，最適輸送距離（Wasserstein 距離）は
地点間のコスト $c(x, y)$ を通じて
\textbf{台の幾何構造}を自然に取り込む．

\begin{example}{離散分布の比較}{discrete-compare}
$\mathcal{X} = \{0, 1, 2\}$ 上の3つの分布を考える：
\begin{align*}
  \mu &= (1,\ 0,\ 0), \\
  \nu_1 &= (0,\ 1,\ 0), \\
  \nu_2 &= (0,\ 0,\ 1).
\end{align*}
KL ダイバージェンスでは $\mathrm{KL}(\mu \| \nu_1)$
も $\mathrm{KL}(\mu \| \nu_2)$ も（台が重ならないため）
ともに $+\infty$ であり区別できない．
しかし，コスト $c(x,y) = |x - y|$ のもとでの
Wasserstein 距離では $W_1(\mu, \nu_1) = 1$ に対し
$W_1(\mu, \nu_2) = 2$ となり，
$\nu_2$ がより遠い分布であるという直感を正しく捉える．
\end{example}

\paragraph{(2) 弱収束と整合する．}
Wasserstein 距離の収束は確率分布の弱収束と整合しており，
生成モデル（GAN 等）の学習で
勾配が消失しにくいという実用上の利点につながる．

\paragraph{(3) 台の幾何的操作が自然に定義できる．}
ユークリッド空間上の内挿・重心計算・勾配といった概念が，
最適輸送の構造を通じて確率分布の空間上に自然に拡張される．
これは画像のモーフィングやドメイン適応などで直接利用される．


% ============================================================
\section{本稿の構成と読み方}\label{sec:outline}

本稿は，Peyré--Cuturi~\cite{PeyreCuturi2019} の全 10 章に対応する
10 章構成である．各章の概要を以下にまとめる．

\medskip

\noindent
\textbf{基礎理論}（第\ref{chap:introduction}--\ref{chap:theoretical}章）
\begin{itemize}
  \item \textbf{第\ref{chap:introduction}章}（本章）：
    歴史的背景と動機，記号の導入．
  \item \textbf{第\ref{chap:theoretical}章}：
    ヒストグラムと測度，Monge 問題，Kantorovich 緩和，双対問題，
    Wasserstein 距離の定義と基本性質．
    本稿全体の数学的基盤を整える．
\end{itemize}

\medskip

\noindent
\textbf{離散問題の計算手法}（第\ref{chap:algorithmic}--\ref{chap:entropic}章）
\begin{itemize}
  \item \textbf{第\ref{chap:algorithmic}章}：
    Kantorovich 問題の線形計画としての解法．
    ネットワーク単体法，双対上昇法，オークションアルゴリズムなど，
    厳密解を求めるアルゴリズムを扱う．
  \item \textbf{第\ref{chap:entropic}章}：
    エントロピー正則化と Sinkhorn アルゴリズム．
    大規模問題に適用可能な近似計算の理論と実装を詳述する．
\end{itemize}

\medskip

\noindent
\textbf{連続と離散の接点}（第\ref{chap:semidiscrete}--\ref{chap:dynamic}章）
\begin{itemize}
  \item \textbf{第\ref{chap:semidiscrete}章}：
    半離散最適輸送．連続分布から離散分布への輸送を扱い，
    $c$変換や確率的最適化手法を紹介する．
  \item \textbf{第\ref{chap:w1}章}：
    $\mathcal{W}_1$ 最適輸送の特殊な構造と，
    距離空間・ユークリッド空間・グラフ上での性質．
  \item \textbf{第\ref{chap:dynamic}章}：
    Benamou--Brenier の動的定式化．
    最適輸送を時間発展する測度の経路として捉える視点を導入する．
\end{itemize}

\medskip

\noindent
\textbf{応用と拡張}（第\ref{chap:divergences}--\ref{chap:extensions}章）
\begin{itemize}
  \item \textbf{第\ref{chap:divergences}章}：
    $\varphi$-ダイバージェンス，MMD など他の統計的距離との関係と比較．
  \item \textbf{第\ref{chap:variational}章}：
    Wasserstein 距離を損失関数として用いる最適化問題．
    重心計算，勾配流，辞書学習などへの応用．
  \item \textbf{第\ref{chap:extensions}章}：
    多周辺問題，非均衡最適輸送，
    スライス Wasserstein 距離，Gromov--Wasserstein 距離など，
    OT の諸拡張．
\end{itemize}

\medskip

各章は原則として独立に読めるように構成しているが，
第\ref{chap:theoretical}章の基礎理論は他の章の前提となるため，
最初に通読することを推奨する．
第\ref{chap:entropic}章の Sinkhorn アルゴリズムは
計算上の応用で繰り返し参照されるので，
応用寄りの興味がある読者は第\ref{chap:theoretical}章のあとに
直接進んでもよい．


% ============================================================
\section{記号一覧}\label{sec:notation}

本稿で用いる主な記号をまとめる．
初出時にも改めて説明するので，
本節は必要に応じて参照すればよい．

\subsection*{基本的な記号}

\begin{tabular}{lp{28em}}
  $[\![ n ]\!]$ & 整数の集合 $\{1, \ldots, n\}$ \\[3pt]
  $\mathbb{1}_{n,m}$ & 成分がすべて $1$ の $n \times m$ 行列．
    $\mathbb{1}_n$ は長さ $n$ の全 1 ベクトル \\[3pt]
  $\mathbb{I}_n$ & $n \times n$ 単位行列 \\[3pt]
  $\diag(u)$ & $u \in \R^n$ を対角成分に持つ対角行列 \\[3pt]
  $\inner{\cdot}{\cdot}$ & ベクトルの内積．
    行列に対しては Frobenius 内積
    $\inner{A}{B} = \tr(A^\top B)$
\end{tabular}

\subsection*{確率分布と測度}

\begin{tabular}{lp{28em}}
  $\Sigma_n$ & $n$ 次元確率単体
    $\bigl\{ \mathbf{a} \in \R^n_+ : \sum_i \mathbf{a}_i = 1 \bigr\}$ \\[3pt]
  $(\mathbf{a}, \mathbf{b})$ & ヒストグラム（離散確率ベクトル），
    $\mathbf{a} \in \Sigma_n,\ \mathbf{b} \in \Sigma_m$ \\[3pt]
  $(\alpha, \beta)$ & 空間 $(\mathcal{X}, \mathcal{Y})$ 上の（一般の）確率測度 \\[3pt]
  $\delta_x$ & 点 $x$ に集中する Dirac 測度 \\[3pt]
  $\alpha = \sum_i \mathbf{a}_i \delta_{x_i}$ & 離散測度の表現 \\[3pt]
  $\rho_\alpha = \frac{\mathrm{d}\alpha}{\mathrm{d}x}$ &
    Lebesgue 測度に関する $\alpha$ の密度関数
\end{tabular}

\subsection*{最適輸送の基本量}

\begin{tabular}{lp{28em}}
  $c(x,y)$ & 地上コスト (ground cost)：$x$ から $y$ への輸送の単位コスト \\[3pt]
  $\mathbf{C}_{i,j}$ & 離散的なコスト行列 $\mathbf{C}_{i,j} = c(x_i, y_j)$ \\[3pt]
  $T : \mathcal{X} \to \mathcal{Y}$ & Monge 写像 \\[3pt]
  $\pi$ & 輸送計画（カップリング測度） \\[3pt]
  $\mathcal{U}(\alpha, \beta)$ &
    $\alpha, \beta$ のカップリング全体の集合．
    離散の場合は $\mathbf{U}(\mathbf{a}, \mathbf{b})$ \\[3pt]
  $(f, g)$ & 双対ポテンシャル．離散の場合は $(\mathbf{f}, \mathbf{g})$ \\[3pt]
  $\mathcal{L}_c(\alpha, \beta)$ &
    コスト $c$ に関する最適輸送コスト．
    離散の場合は $\mathbf{L_C}(\mathbf{a}, \mathbf{b})$ \\[3pt]
  $\mathcal{W}_p(\alpha, \beta)$ &
    $p$-Wasserstein 距離．
    離散の場合は $\mathbf{W}_p(\mathbf{a}, \mathbf{b})$
\end{tabular}

\subsection*{エントロピー正則化}

\begin{tabular}{lp{28em}}
  $\varepsilon$ & 正則化パラメータ \\[3pt]
  $\mathbf{K} \defeq e^{-\mathbf{C}/\varepsilon}$ &
    Gibbs カーネル（Sinkhorn アルゴリズムの核） \\[3pt]
  $(\mathbf{u}, \mathbf{v})$ &
    Sinkhorn スケーリング変数
    $(\mathbf{u}, \mathbf{v}) \defeq (e^{\mathbf{f}/\varepsilon},\, e^{\mathbf{g}/\varepsilon})$
\end{tabular}

\subsection*{その他の演算}

\begin{tabular}{lp{28em}}
  $f \oplus g$ & $(f \oplus g)(x,y) \defeq f(x) + g(y)$．
    離散の場合 $\mathbf{f} \oplus \mathbf{g}
    = \mathbf{f}\, \mathbb{1}_m^\top + \mathbb{1}_n\, \mathbf{g}^\top$ \\[3pt]
  $\alpha \otimes \beta$ & 積測度：
    $\int g\, \mathrm{d}(\alpha \otimes \beta)
    = \int\!\!\int g(x,y)\, \mathrm{d}\alpha(x)\, \mathrm{d}\beta(y)$ \\[3pt]
  $\mathbf{u} \odot \mathbf{v}$ & 成分ごとの積（Hadamard 積）
\end{tabular}
